
\documentclass{article}
\usepackage{colortbl}
\usepackage{makecell}
\usepackage{multirow}
\usepackage{supertabular}

\begin{document}

\newcounter{utterance}

\centering \large Interaction Transcript for game `clean\_up', experiment `0\_easy\_5obj\_ru', episode 3 with gemma{-}4{-}26b{-}a4b{-}it.
\vspace{24pt}

{ \footnotesize  \setcounter{utterance}{1}
\setlength{\tabcolsep}{0pt}
\begin{supertabular}{c@{$\;$}|p{.15\linewidth}@{}p{.15\linewidth}p{.15\linewidth}p{.15\linewidth}p{.15\linewidth}p{.15\linewidth}}
    \# & $\;$A & \multicolumn{4}{c}{Game Master} & $\;\:$B\\
    \hline

    \theutterance \stepcounter{utterance}  
    & & \multicolumn{4}{p{0.6\linewidth}}{
        \cellcolor[rgb]{0.9,0.9,0.9}{
            \makecell[{{p{\linewidth}}}]{
                \texttt{\tiny{[P1$\langle$GM]}}
                \texttt{Я твой ведущий, и ты играешь в совместную игру, где игровым полем служит следующая сетка:} \\
\\ 
\texttt{```} \\
\texttt{ 1234567} \\
\texttt{╔════╤══╗} \\
\texttt{║H◌◌◌│◌◌║ 1} \\
\texttt{║◌W◌◌│◌◌║ 2} \\
\texttt{║◌C◌◌│◌◌║ 3} \\
\texttt{║T◌◌◌│I◌║ 4} \\
\texttt{╟──┐◌│◌◌║ 5} \\
\texttt{║◌◌├─┼──╢ 6} \\
\texttt{║◌◌│◌│◌◌║ 7} \\
\texttt{╚══╧═╧══╝} \\
\\ 
\texttt{```} \\
\\ 
\texttt{* По верхнему краю показаны x{-}координаты, увеличивающиеся вправо, а по правому краю — y{-}координаты, увеличивающиеся вниз.} \\
\texttt{* На твоей сетке случайным образом размещены следующие объекты: 'W', 'I', 'T', 'C', 'H'.} \\
\\ 
\\ 
\texttt{Другой игрок видит вариант игрового поля, где объекты размещены в других случайных местах. Ты не видишь поле другого игрока, и он не видит твоё.} \\
\\ 
\\ 
\texttt{**Цель:**} \\
\\ 
\texttt{Оба игрока должны перемещать объекты на своём фоне так, чтобы одинаковые объекты оказались в одних и тех же координатах. Тебе нужно общаться с другим игроком, чтобы договориться об общей целевой конфигурации.} \\
\\ 
\\ 
\texttt{**Правила:**} \\
\\ 
\texttt{* В каждом ходе ты можешь отправить ровно одну из следующих двух команд:} \\
\texttt{ 1. `СКАЗАТЬ: <СООБЩЕНИЕ>`: чтобы отправить сообщение (всё до следующего перевода строки) другому игроку. Я перешлю его твоему партнёру.} \\
\texttt{2. `ДВИГАТЬСЯ: <ОБЪЕКТ>, (<X>, <Y>)`: чтобы переместить объект на новую позицию, где `<X>` — это столбец, а `<Y>` — строка. Я сообщу тебе, был ли ход успешным или нет.} \\
\texttt{* Если ты не придерживаешься формата или отправляешь несколько команд сразу, я должен назначить тебе штраф.} \\
\texttt{* Если оба игрока наберут больше чем 12 штрафов, вы оба проиграете.} \\
\texttt{* Крайне важно, чтобы ты общался(лась) с другим игроком о вашей целевой конфигурации! *Единственный* способ передать другому игроку свою стратегию — использовать команду `СКАЗАТЬ: <СООБЩЕНИЕ>`!} \\
\\ 
\\ 
\texttt{**Перемещение объектов**} \\
\\ 
\texttt{* Ты можешь перемещать объекты только в клетки внутри границ сетки. Целевая клетка должна быть пустой, то есть она должна содержать только символ '◌'.} \\
\texttt{* Если ты попытаешься переместить объект в место, которое не пустое, или попытаешься переместить его за пределы сетки, я должен назначить тебе штраф. У тебя будет ещё одна попытка.} \\
\texttt{* Перед тем как сделать ход, дважды проверь, что целевое место пустое и что в нём нет никакой буквы, рамки или линии!} \\
\\ 
\\ 
\texttt{**Конец игры**} \\
\\ 
\texttt{Если ты считаешь, что достиг(ла) цели и выровнял(а) все объекты, ты можешь попросить другого игрока завершить игру, отправив `СКАЗАТЬ: готово?`. Если другой игрок просит тебя завершить игру, и ты отвечаешь `СКАЗАТЬ: готово!`, игра завершится.} \\
\\ 
\texttt{Оба игрока выигрывают, если игра закончится в пределах 20 раундов, где один раунд определяется как то, что оба игрока отправили по одной корректной команде.} \\
\\ 
\texttt{**Подсчёт очков:**} \\
\\ 
\texttt{Чем ближе одинаковые объекты расположены на обоих игровых полях, тем больше очков ты получаешь. Штрафы уменьшают твои очки. Сможешь побить рекорд?} \\
\\ 
\\ 
\texttt{Пожалуйста, отправь сообщение другому игроку, чтобы начать игру!} \\
            }
        }
    }
    & \\ \\

    \theutterance \stepcounter{utterance}  
    & \multicolumn{4}{p{0.6\linewidth}}{
        \cellcolor[rgb]{0.9,0.9,0.9}{
            \makecell[{{p{\linewidth}}}]{
                \texttt{\tiny{[P1$\rangle$GM]}}
                \texttt{СКАЗАТЬ: Привет! Давай договоримся о целевой конфигурации. Мои объекты сейчас на: H(1,1), W(2,2), C(2,3), T(1,4), I(5,4). Предлагаю собрать их в ряд или в компактную группу. Какие у тебя координаты объектов?} \\
            }
        }
    }
    & & \\ \\

    \theutterance \stepcounter{utterance}  
    & & & \multicolumn{2}{p{0.3\linewidth}}{
        \cellcolor[rgb]{0.9,0.9,0.9}{
            \makecell[{{p{\linewidth}}}]{
                \texttt{\tiny{[GM$|$GM]}}
                \texttt{СКАЗАТЬ: Привет! Давай договоримся о целевой конфигурации. Мои объекты сейчас на: H(1,1), W(2,2), C(2,3), T(1,4), I(5,4). Предлагаю собрать их в ряд или в компактную группу. Какие у тебя координаты объектов?} \\
            }
        }
    }
    & & \\ \\

    \theutterance \stepcounter{utterance}  
    & & \multicolumn{4}{p{0.6\linewidth}}{
        \cellcolor[rgb]{0.9,0.9,0.9}{
            \makecell[{{p{\linewidth}}}]{
                \texttt{\tiny{[GM$\rangle$P2]}}
                \texttt{Я твой ведущий, и ты играешь в совместную игру, где игровым полем служит следующая сетка:} \\
\\ 
\texttt{```} \\
\texttt{ 1234567} \\
\texttt{╔════╤══╗} \\
\texttt{║◌◌◌T│◌◌║ 1} \\
\texttt{║◌◌◌◌│◌◌║ 2} \\
\texttt{║◌◌◌◌│◌◌║ 3} \\
\texttt{║W◌◌◌│◌◌║ 4} \\
\texttt{╟──┐◌│◌◌║ 5} \\
\texttt{║H◌├─┼──╢ 6} \\
\texttt{║I◌│◌│◌C║ 7} \\
\texttt{╚══╧═╧══╝} \\
\\ 
\texttt{```} \\
\\ 
\texttt{* По верхнему краю показаны x{-}координаты, увеличивающиеся вправо, а по правому краю — y{-}координаты, увеличивающиеся вниз.} \\
\texttt{* На твоей сетке случайным образом размещены следующие объекты: 'W', 'I', 'T', 'C', 'H'.} \\
\\ 
\\ 
\texttt{Другой игрок видит вариант игрового поля, где объекты размещены в других случайных местах. Ты не видишь поле другого игрока, и он не видит твоё.} \\
\\ 
\\ 
\texttt{**Цель:**} \\
\\ 
\texttt{Оба игрока должны перемещать объекты на своём фоне так, чтобы одинаковые объекты оказались в одних и тех же координатах. Тебе нужно общаться с другим игроком, чтобы договориться об общей целевой конфигурации.} \\
\\ 
\\ 
\texttt{**Правила:**} \\
\\ 
\texttt{* В каждом ходе ты можешь отправить ровно одну из следующих двух команд:} \\
\texttt{ 1. `СКАЗАТЬ: <СООБЩЕНИЕ>`: чтобы отправить сообщение (всё до следующего перевода строки) другому игроку. Я перешлю его твоему партнёру.} \\
\texttt{2. `ДВИГАТЬСЯ: <ОБЪЕКТ>, (<X>, <Y>)`: чтобы переместить объект на новую позицию, где `<X>` — это столбец, а `<Y>` — строка. Я сообщу тебе, был ли ход успешным или нет.} \\
\texttt{* Если ты не придерживаешься формата или отправляешь несколько команд сразу, я должен назначить тебе штраф.} \\
\texttt{* Если оба игрока наберут больше чем 12 штрафов, вы оба проиграете.} \\
\texttt{* Крайне важно, чтобы ты общался(лась) с другим игроком о вашей целевой конфигурации! *Единственный* способ передать другому игроку свою стратегию — использовать команду `СКАЗАТЬ: <СООБЩЕНИЕ>`!} \\
\\ 
\\ 
\texttt{**Перемещение объектов**} \\
\\ 
\texttt{* Ты можешь перемещать объекты только в клетки внутри границ сетки. Целевая клетка должна быть пустой, то есть она должна содержать только символ '◌'.} \\
\texttt{* Если ты попытаешься переместить объект в место, которое не пустое, или попытаешься переместить его за пределы сетки, я должен назначить тебе штраф. У тебя будет ещё одна попытка.} \\
\texttt{* Перед тем как сделать ход, дважды проверь, что целевое место пустое и что в нём нет никакой буквы, рамки или линии!} \\
\\ 
\\ 
\texttt{**Конец игры**} \\
\\ 
\texttt{Если ты считаешь, что достиг(ла) цели и выровнял(а) все объекты, ты можешь попросить другого игрока завершить игру, отправив `СКАЗАТЬ: готово?`. Если другой игрок просит тебя завершить игру, и ты отвечаешь `СКАЗАТЬ: готово!`, игра завершится.} \\
\\ 
\texttt{Оба игрока выигрывают, если игра закончится в пределах 20 раундов, где один раунд определяется как то, что оба игрока отправили по одной корректной команде.} \\
\\ 
\texttt{**Подсчёт очков:**} \\
\\ 
\texttt{Чем ближе одинаковые объекты расположены на обоих игровых полях, тем больше очков ты получаешь. Штрафы уменьшают твои очки. Сможешь побить рекорд?} \\
\\ 
\texttt{Другой игрок начал игру, отправив это сообщение:} \\
\\ 
\texttt{"Привет! Давай договоримся о целевой конфигурации. Мои объекты сейчас на: H(1,1), W(2,2), C(2,3), T(1,4), I(5,4). Предлагаю собрать их в ряд или в компактную группу. Какие у тебя координаты объектов?"} \\
\\ 
\texttt{Какая твоя первая команда?} \\
            }
        }
    }
    & \\ \\

    \theutterance \stepcounter{utterance}  
    & & & \multicolumn{4}{p{0.6\linewidth}}{
        \cellcolor[rgb]{0.9,0.9,0.9}{
            \makecell[{{p{\linewidth}}}]{
                \texttt{\tiny{[GM$\langle$P2]}}
                \texttt{СКАЗАТЬ: Привет! Мои координаты: T(4,1), W(1,4), H(1,6), I(1,7), C(6,7). Давай попробуем собрать их в компактную группу в центре или в левой части поля. Предлагаю такую цель: T(2,2), W(2,3), H(2,4), I(2,5), C(2,6). Что думаешь?} \\
            }
        }
    }
     \\ \\

    \theutterance \stepcounter{utterance}  
    & & & \multicolumn{2}{p{0.3\linewidth}}{
        \cellcolor[rgb]{0.9,0.9,0.9}{
            \makecell[{{p{\linewidth}}}]{
                \texttt{\tiny{[GM$|$GM]}}
                \texttt{СКАЗАТЬ: Привет! Мои координаты: T(4,1), W(1,4), H(1,6), I(1,7), C(6,7). Давай попробуем собрать их в компактную группу в центре или в левой части поля. Предлагаю такую цель: T(2,2), W(2,3), H(2,4), I(2,5), C(2,6). Что думаешь?} \\
            }
        }
    }
    & & \\ \\

    \theutterance \stepcounter{utterance}  
    & & \multicolumn{4}{p{0.6\linewidth}}{
        \cellcolor[rgb]{0.9,0.9,0.9}{
            \makecell[{{p{\linewidth}}}]{
                \texttt{\tiny{[P1$\langle$GM]}}
                \texttt{Ваше сообщение было передано другому игроку.} \\
\\ 
\\ 
\texttt{Сейчас вы играете раунд 1 из максимум 20.} \\
\texttt{Вы вместе накопили 0 из 12 штрафов.} \\
\texttt{Другой игрок отправил это сообщение:} \\
\\ 
\\ 
\texttt{"Привет! Мои координаты: T(4,1), W(1,4), H(1,6), I(1,7), C(6,7). Давай попробуем собрать их в компактную группу в центре или в левой части поля. Предлагаю такую цель: T(2,2), W(2,3), H(2,4), I(2,5), C(2,6). Что думаешь?"} \\
\\ 
\\ 
\texttt{Какая твоя следующая команда?} \\
            }
        }
    }
    & \\ \\

    \theutterance \stepcounter{utterance}  
    & \multicolumn{4}{p{0.6\linewidth}}{
        \cellcolor[rgb]{0.9,0.9,0.9}{
            \makecell[{{p{\linewidth}}}]{
                \texttt{\tiny{[P1$\rangle$GM]}}
                \texttt{СКАЗАТЬ: Отличный план! Твои координаты цели подходят. Давай придерживаться твоего варианта: T(2,2), W(2,3), H(2,4), I(2,5), C(2,6). Я начну двигать свои объекты к этим позициям. Мой первый ход: ДВИГАТЬСЯ: H, (2,4).} \\
            }
        }
    }
    & & \\ \\

    \theutterance \stepcounter{utterance}  
    & & & \multicolumn{2}{p{0.3\linewidth}}{
        \cellcolor[rgb]{0.9,0.9,0.9}{
            \makecell[{{p{\linewidth}}}]{
                \texttt{\tiny{[GM$|$GM]}}
                \texttt{СКАЗАТЬ: Отличный план! Твои координаты цели подходят. Давай придерживаться твоего варианта: T(2,2), W(2,3), H(2,4), I(2,5), C(2,6). Я начну двигать свои объекты к этим позициям. Мой первый ход: ДВИГАТЬСЯ: H, (2,4).} \\
            }
        }
    }
    & & \\ \\

    \theutterance \stepcounter{utterance}  
    & & & \multicolumn{2}{p{0.3\linewidth}}{
        \cellcolor[rgb]{0.9,0.9,0.9}{
            \makecell[{{p{\linewidth}}}]{
                \texttt{\tiny{[GM$|$GM]}}
                \texttt{Invalid response format: СКАЗАТЬ: Отличный план! Твои координаты цели подходят. Давай придерживаться твоего варианта: T(2,2), W(2,3), H(2,4), I(2,5), C(2,6). Я начну двигать свои объекты к этим позициям. Мой первый ход: ДВИГАТЬСЯ: H, (2,4).} \\
            }
        }
    }
    & & \\ \\

    \theutterance \stepcounter{utterance}  
    & & \multicolumn{4}{p{0.6\linewidth}}{
        \cellcolor[rgb]{0.9,0.9,0.9}{
            \makecell[{{p{\linewidth}}}]{
                \texttt{\tiny{[P1$\langle$GM]}}
                \texttt{Штраф: Ваше сообщение содержит более одной команды!} \\
\\ 
\texttt{Убедитесь, что ваш ответ содержит только `СКАЗАТЬ: <СООБЩЕНИЕ>` или `ДВИГАТЬСЯ: <ОБЪЕКТ>, (<X>, <Y>)` и ничего больше!} \\
\\ 
\texttt{Вы вместе накопили 1 из 12 штрафов.} \\
\texttt{Пожалуйста, попробуйте ещё раз!} \\
            }
        }
    }
    & \\ \\

    \theutterance \stepcounter{utterance}  
    & \multicolumn{4}{p{0.6\linewidth}}{
        \cellcolor[rgb]{0.9,0.9,0.9}{
            \makecell[{{p{\linewidth}}}]{
                \texttt{\tiny{[P1$\rangle$GM]}}
                \texttt{ДВИГАТЬСЯ: H, (2,4)} \\
            }
        }
    }
    & & \\ \\

    \theutterance \stepcounter{utterance}  
    & & & \multicolumn{2}{p{0.3\linewidth}}{
        \cellcolor[rgb]{0.9,0.9,0.9}{
            \makecell[{{p{\linewidth}}}]{
                \texttt{\tiny{[GM$|$GM]}}
                \texttt{ДВИГАТЬСЯ: H, (2,4)} \\
            }
        }
    }
    & & \\ \\

    \theutterance \stepcounter{utterance}  
    & & & \multicolumn{2}{p{0.3\linewidth}}{
        \cellcolor[rgb]{0.9,0.9,0.9}{
            \makecell[{{p{\linewidth}}}]{
                \texttt{\tiny{[GM$|$GM]}}
                \texttt{H moved to (2, 4)} \\
            }
        }
    }
    & & \\ \\

    \theutterance \stepcounter{utterance}  
    & & \multicolumn{4}{p{0.6\linewidth}}{
        \cellcolor[rgb]{0.9,0.9,0.9}{
            \makecell[{{p{\linewidth}}}]{
                \texttt{\tiny{[GM$\rangle$P2]}}
                \texttt{Ваше сообщение было передано другому игроку.} \\
\\ 
\\ 
\texttt{Сейчас вы играете раунд 2 из максимум 20.} \\
\texttt{Вы вместе накопили 1 из 12 штрафов.} \\
\texttt{Другой игрок переместил объект на своей сетке.} \\
\\ 
\\ 
\texttt{Какая твоя следующая команда?} \\
            }
        }
    }
    & \\ \\

    \theutterance \stepcounter{utterance}  
    & & & \multicolumn{4}{p{0.6\linewidth}}{
        \cellcolor[rgb]{0.9,0.9,0.9}{
            \makecell[{{p{\linewidth}}}]{
                \texttt{\tiny{[GM$\langle$P2]}}
                \texttt{Чтобы я мог эффективно продолжать игру, мне нужно знать, какой ход сделал другой игрок (какой объект он переместил и куда) и ответил ли он на моё предложение по координатам.} \\
\\ 
\texttt{Пожалуйста, сообщи мне сообщение или действие другого игрока.} \\
            }
        }
    }
     \\ \\

    \theutterance \stepcounter{utterance}  
    & & & \multicolumn{2}{p{0.3\linewidth}}{
        \cellcolor[rgb]{0.9,0.9,0.9}{
            \makecell[{{p{\linewidth}}}]{
                \texttt{\tiny{[GM$|$GM]}}
                \texttt{Чтобы я мог эффективно продолжать игру, мне нужно знать, какой ход сделал другой игрок (какой объект он переместил и куда) и ответил ли он на моё предложение по координатам.} \\
\\ 
\texttt{Пожалуйста, сообщи мне сообщение или действие другого игрока.} \\
            }
        }
    }
    & & \\ \\

    \theutterance \stepcounter{utterance}  
    & & & \multicolumn{2}{p{0.3\linewidth}}{
        \cellcolor[rgb]{0.9,0.9,0.9}{
            \makecell[{{p{\linewidth}}}]{
                \texttt{\tiny{[GM$|$GM]}}
                \texttt{Invalid response format} \\
            }
        }
    }
    & & \\ \\

    \theutterance \stepcounter{utterance}  
    & & \multicolumn{4}{p{0.6\linewidth}}{
        \cellcolor[rgb]{0.9,0.9,0.9}{
            \makecell[{{p{\linewidth}}}]{
                \texttt{\tiny{[GM$\rangle$P2]}}
                \texttt{Штраф: Ваше сообщение не соответствует ожидаемому формату!} \\
\\ 
\texttt{Убедитесь, что ваш ответ содержит только `СКАЗАТЬ: <СООБЩЕНИЕ>` или `ДВИГАТЬСЯ: <ОБЪЕКТ>, (<X>, <Y>)` и ничего больше!} \\
\\ 
\texttt{Вы вместе накопили 2 из 12 штрафов.} \\
\texttt{Пожалуйста, попробуйте ещё раз!} \\
            }
        }
    }
    & \\ \\

    \theutterance \stepcounter{utterance}  
    & & & \multicolumn{4}{p{0.6\linewidth}}{
        \cellcolor[rgb]{0.9,0.9,0.9}{
            \makecell[{{p{\linewidth}}}]{
                \texttt{\tiny{[GM$\langle$P2]}}
                \texttt{СКАЗАТЬ: Извини за ошибку. Жду твоего ответа по поводу предложенных координат: T(2,2), W(2,3), H(2,4), I(2,5), C(2,6).} \\
            }
        }
    }
     \\ \\

    \theutterance \stepcounter{utterance}  
    & & & \multicolumn{2}{p{0.3\linewidth}}{
        \cellcolor[rgb]{0.9,0.9,0.9}{
            \makecell[{{p{\linewidth}}}]{
                \texttt{\tiny{[GM$|$GM]}}
                \texttt{СКАЗАТЬ: Извини за ошибку. Жду твоего ответа по поводу предложенных координат: T(2,2), W(2,3), H(2,4), I(2,5), C(2,6).} \\
            }
        }
    }
    & & \\ \\

    \theutterance \stepcounter{utterance}  
    & & \multicolumn{4}{p{0.6\linewidth}}{
        \cellcolor[rgb]{0.9,0.9,0.9}{
            \makecell[{{p{\linewidth}}}]{
                \texttt{\tiny{[P1$\langle$GM]}}
                \texttt{Успешно перемещён объект "H" в (2,4). Ваша обновлённая сетка выглядит так:} \\
\texttt{```} \\
\texttt{ 1234567} \\
\texttt{╔════╤══╗} \\
\texttt{║◌◌◌◌│◌◌║ 1} \\
\texttt{║◌W◌◌│◌◌║ 2} \\
\texttt{║◌C◌◌│◌◌║ 3} \\
\texttt{║TH◌◌│I◌║ 4} \\
\texttt{╟──┐◌│◌◌║ 5} \\
\texttt{║◌◌├─┼──╢ 6} \\
\texttt{║◌◌│◌│◌◌║ 7} \\
\texttt{╚══╧═╧══╝} \\
\texttt{```} \\
\texttt{Сейчас вы играете раунд 2 из максимум 20.} \\
\texttt{Вы вместе накопили 2 из 12 штрафов.} \\
\texttt{Другой игрок отправил это сообщение:} \\
\\ 
\\ 
\texttt{"Извини за ошибку. Жду твоего ответа по поводу предложенных координат: T(2,2), W(2,3), H(2,4), I(2,5), C(2,6)."} \\
\\ 
\\ 
\texttt{Какая твоя следующая команда?} \\
            }
        }
    }
    & \\ \\

    \theutterance \stepcounter{utterance}  
    & \multicolumn{4}{p{0.6\linewidth}}{
        \cellcolor[rgb]{0.9,0.9,0.9}{
            \makecell[{{p{\linewidth}}}]{
                \texttt{\tiny{[P1$\rangle$GM]}}
                \texttt{ДВИГАТЬСЯ: W, (2,3)} \\
            }
        }
    }
    & & \\ \\

    \theutterance \stepcounter{utterance}  
    & & & \multicolumn{2}{p{0.3\linewidth}}{
        \cellcolor[rgb]{0.9,0.9,0.9}{
            \makecell[{{p{\linewidth}}}]{
                \texttt{\tiny{[GM$|$GM]}}
                \texttt{ДВИГАТЬСЯ: W, (2,3)} \\
            }
        }
    }
    & & \\ \\

    \theutterance \stepcounter{utterance}  
    & & & \multicolumn{2}{p{0.3\linewidth}}{
        \cellcolor[rgb]{0.9,0.9,0.9}{
            \makecell[{{p{\linewidth}}}]{
                \texttt{\tiny{[GM$|$GM]}}
                \texttt{Штраф: (2,3) не пусто, а содержит "C".} \\
\\ 
\texttt{Вы вместе накопили 3 из 12 штрафов.} \\
\texttt{Пожалуйста, попробуйте ещё раз!} \\
            }
        }
    }
    & & \\ \\

    \theutterance \stepcounter{utterance}  
    & & \multicolumn{4}{p{0.6\linewidth}}{
        \cellcolor[rgb]{0.9,0.9,0.9}{
            \makecell[{{p{\linewidth}}}]{
                \texttt{\tiny{[P1$\langle$GM]}}
                \texttt{Штраф: (2,3) не пусто, а содержит "C".} \\
\\ 
\texttt{Вы вместе накопили 3 из 12 штрафов.} \\
\texttt{Пожалуйста, попробуйте ещё раз!} \\
            }
        }
    }
    & \\ \\

    \theutterance \stepcounter{utterance}  
    & \multicolumn{4}{p{0.6\linewidth}}{
        \cellcolor[rgb]{0.9,0.9,0.9}{
            \makecell[{{p{\linewidth}}}]{
                \texttt{\tiny{[P1$\rangle$GM]}}
                \texttt{ДВИГАТЬСЯ: C, (2,6)} \\
            }
        }
    }
    & & \\ \\

    \theutterance \stepcounter{utterance}  
    & & & \multicolumn{2}{p{0.3\linewidth}}{
        \cellcolor[rgb]{0.9,0.9,0.9}{
            \makecell[{{p{\linewidth}}}]{
                \texttt{\tiny{[GM$|$GM]}}
                \texttt{ДВИГАТЬСЯ: C, (2,6)} \\
            }
        }
    }
    & & \\ \\

    \theutterance \stepcounter{utterance}  
    & & & \multicolumn{2}{p{0.3\linewidth}}{
        \cellcolor[rgb]{0.9,0.9,0.9}{
            \makecell[{{p{\linewidth}}}]{
                \texttt{\tiny{[GM$|$GM]}}
                \texttt{C moved to (2, 6)} \\
            }
        }
    }
    & & \\ \\

    \theutterance \stepcounter{utterance}  
    & & \multicolumn{4}{p{0.6\linewidth}}{
        \cellcolor[rgb]{0.9,0.9,0.9}{
            \makecell[{{p{\linewidth}}}]{
                \texttt{\tiny{[GM$\rangle$P2]}}
                \texttt{Ваше сообщение было передано другому игроку.} \\
\\ 
\\ 
\texttt{Сейчас вы играете раунд 3 из максимум 20.} \\
\texttt{Вы вместе накопили 3 из 12 штрафов.} \\
\texttt{Другой игрок переместил объект на своей сетке.} \\
\\ 
\\ 
\texttt{Какая твоя следующая команда?} \\
            }
        }
    }
    & \\ \\

    \theutterance \stepcounter{utterance}  
    & & & \multicolumn{4}{p{0.6\linewidth}}{
        \cellcolor[rgb]{0.9,0.9,0.9}{
            \makecell[{{p{\linewidth}}}]{
                \texttt{\tiny{[GM$\langle$P2]}}
                \texttt{Мне нужно знать, что сделал другой игрок (какой объект он переместил и куда), чтобы я мог продолжить игру. Пожалуйста, сообщи мне его ход.} \\
            }
        }
    }
     \\ \\

    \theutterance \stepcounter{utterance}  
    & & & \multicolumn{2}{p{0.3\linewidth}}{
        \cellcolor[rgb]{0.9,0.9,0.9}{
            \makecell[{{p{\linewidth}}}]{
                \texttt{\tiny{[GM$|$GM]}}
                \texttt{Мне нужно знать, что сделал другой игрок (какой объект он переместил и куда), чтобы я мог продолжить игру. Пожалуйста, сообщи мне его ход.} \\
            }
        }
    }
    & & \\ \\

    \theutterance \stepcounter{utterance}  
    & & & \multicolumn{2}{p{0.3\linewidth}}{
        \cellcolor[rgb]{0.9,0.9,0.9}{
            \makecell[{{p{\linewidth}}}]{
                \texttt{\tiny{[GM$|$GM]}}
                \texttt{Invalid response format} \\
            }
        }
    }
    & & \\ \\

    \theutterance \stepcounter{utterance}  
    & & \multicolumn{4}{p{0.6\linewidth}}{
        \cellcolor[rgb]{0.9,0.9,0.9}{
            \makecell[{{p{\linewidth}}}]{
                \texttt{\tiny{[GM$\rangle$P2]}}
                \texttt{Штраф: Ваше сообщение не соответствует ожидаемому формату!} \\
\\ 
\texttt{Убедитесь, что ваш ответ содержит только `СКАЗАТЬ: <СООБЩЕНИЕ>` или `ДВИГАТЬСЯ: <ОБЪЕКТ>, (<X>, <Y>)` и ничего больше!} \\
\\ 
\texttt{Вы вместе накопили 4 из 12 штрафов.} \\
\texttt{Пожалуйста, попробуйте ещё раз!} \\
            }
        }
    }
    & \\ \\

    \theutterance \stepcounter{utterance}  
    & & & \multicolumn{4}{p{0.6\linewidth}}{
        \cellcolor[rgb]{0.9,0.9,0.9}{
            \makecell[{{p{\linewidth}}}]{
                \texttt{\tiny{[GM$\langle$P2]}}
                \texttt{СКАЗАТЬ: Пожалуйста, сообщи, согласен ли ты на мои координаты (T:2,2, W:2,3, H:2,4, I:2,5, C:2,6) или предложи свои. Также напиши, какой объект ты переместил.} \\
            }
        }
    }
     \\ \\

    \theutterance \stepcounter{utterance}  
    & & & \multicolumn{2}{p{0.3\linewidth}}{
        \cellcolor[rgb]{0.9,0.9,0.9}{
            \makecell[{{p{\linewidth}}}]{
                \texttt{\tiny{[GM$|$GM]}}
                \texttt{СКАЗАТЬ: Пожалуйста, сообщи, согласен ли ты на мои координаты (T:2,2, W:2,3, H:2,4, I:2,5, C:2,6) или предложи свои. Также напиши, какой объект ты переместил.} \\
            }
        }
    }
    & & \\ \\

    \theutterance \stepcounter{utterance}  
    & & \multicolumn{4}{p{0.6\linewidth}}{
        \cellcolor[rgb]{0.9,0.9,0.9}{
            \makecell[{{p{\linewidth}}}]{
                \texttt{\tiny{[P1$\langle$GM]}}
                \texttt{Успешно перемещён объект "C" в (2,6). Ваша обновлённая сетка выглядит так:} \\
\texttt{```} \\
\texttt{ 1234567} \\
\texttt{╔════╤══╗} \\
\texttt{║◌◌◌◌│◌◌║ 1} \\
\texttt{║◌W◌◌│◌◌║ 2} \\
\texttt{║◌◌◌◌│◌◌║ 3} \\
\texttt{║TH◌◌│I◌║ 4} \\
\texttt{╟──┐◌│◌◌║ 5} \\
\texttt{║◌C├─┼──╢ 6} \\
\texttt{║◌◌│◌│◌◌║ 7} \\
\texttt{╚══╧═╧══╝} \\
\texttt{```} \\
\texttt{Сейчас вы играете раунд 3 из максимум 20.} \\
\texttt{Вы вместе накопили 4 из 12 штрафов.} \\
\texttt{Другой игрок отправил это сообщение:} \\
\\ 
\\ 
\texttt{"Пожалуйста, сообщи, согласен ли ты на мои координаты (T:2,2, W:2,3, H:2,4, I:2,5, C:2,6) или предложи свои. Также напиши, какой объект ты переместил."} \\
\\ 
\\ 
\texttt{Какая твоя следующая команда?} \\
            }
        }
    }
    & \\ \\

    \theutterance \stepcounter{utterance}  
    & \multicolumn{4}{p{0.6\linewidth}}{
        \cellcolor[rgb]{0.9,0.9,0.9}{
            \makecell[{{p{\linewidth}}}]{
                \texttt{\tiny{[P1$\rangle$GM]}}
                \texttt{СКАЗАТЬ: Да, я согласен на твои координаты (T:2,2, W:2,3, H:2,4, I:2,5, C:2,6). Я только что переместил C в (2,6).} \\
            }
        }
    }
    & & \\ \\

    \theutterance \stepcounter{utterance}  
    & & & \multicolumn{2}{p{0.3\linewidth}}{
        \cellcolor[rgb]{0.9,0.9,0.9}{
            \makecell[{{p{\linewidth}}}]{
                \texttt{\tiny{[GM$|$GM]}}
                \texttt{СКАЗАТЬ: Да, я согласен на твои координаты (T:2,2, W:2,3, H:2,4, I:2,5, C:2,6). Я только что переместил C в (2,6).} \\
            }
        }
    }
    & & \\ \\

    \theutterance \stepcounter{utterance}  
    & & \multicolumn{4}{p{0.6\linewidth}}{
        \cellcolor[rgb]{0.9,0.9,0.9}{
            \makecell[{{p{\linewidth}}}]{
                \texttt{\tiny{[GM$\rangle$P2]}}
                \texttt{Ваше сообщение было передано другому игроку.} \\
\\ 
\\ 
\texttt{Сейчас вы играете раунд 4 из максимум 20.} \\
\texttt{Вы вместе накопили 4 из 12 штрафов.} \\
\texttt{Другой игрок отправил это сообщение:} \\
\\ 
\\ 
\texttt{"Да, я согласен на твои координаты (T:2,2, W:2,3, H:2,4, I:2,5, C:2,6). Я только что переместил C в (2,6)."} \\
\\ 
\\ 
\texttt{Какая твоя следующая команда?} \\
            }
        }
    }
    & \\ \\

    \theutterance \stepcounter{utterance}  
    & & & \multicolumn{4}{p{0.6\linewidth}}{
        \cellcolor[rgb]{0.9,0.9,0.9}{
            \makecell[{{p{\linewidth}}}]{
                \texttt{\tiny{[GM$\langle$P2]}}
                \texttt{ДВИГАТЬСЯ: C, (2, 6)} \\
            }
        }
    }
     \\ \\

    \theutterance \stepcounter{utterance}  
    & & & \multicolumn{2}{p{0.3\linewidth}}{
        \cellcolor[rgb]{0.9,0.9,0.9}{
            \makecell[{{p{\linewidth}}}]{
                \texttt{\tiny{[GM$|$GM]}}
                \texttt{ДВИГАТЬСЯ: C, (2, 6)} \\
            }
        }
    }
    & & \\ \\

    \theutterance \stepcounter{utterance}  
    & & & \multicolumn{2}{p{0.3\linewidth}}{
        \cellcolor[rgb]{0.9,0.9,0.9}{
            \makecell[{{p{\linewidth}}}]{
                \texttt{\tiny{[GM$|$GM]}}
                \texttt{C moved to (2, 6)} \\
            }
        }
    }
    & & \\ \\

    \theutterance \stepcounter{utterance}  
    & & \multicolumn{4}{p{0.6\linewidth}}{
        \cellcolor[rgb]{0.9,0.9,0.9}{
            \makecell[{{p{\linewidth}}}]{
                \texttt{\tiny{[P1$\langle$GM]}}
                \texttt{Ваше сообщение было передано другому игроку.} \\
\\ 
\\ 
\texttt{Сейчас вы играете раунд 4 из максимум 20.} \\
\texttt{Вы вместе накопили 4 из 12 штрафов.} \\
\texttt{Другой игрок переместил объект на своей сетке.} \\
\\ 
\\ 
\texttt{Какая твоя следующая команда?} \\
            }
        }
    }
    & \\ \\

    \theutterance \stepcounter{utterance}  
    & \multicolumn{4}{p{0.6\linewidth}}{
        \cellcolor[rgb]{0.9,0.9,0.9}{
            \makecell[{{p{\linewidth}}}]{
                \texttt{\tiny{[P1$\rangle$GM]}}
                \texttt{ДВИГАТЬСЯ: W, (2,3)} \\
            }
        }
    }
    & & \\ \\

    \theutterance \stepcounter{utterance}  
    & & & \multicolumn{2}{p{0.3\linewidth}}{
        \cellcolor[rgb]{0.9,0.9,0.9}{
            \makecell[{{p{\linewidth}}}]{
                \texttt{\tiny{[GM$|$GM]}}
                \texttt{ДВИГАТЬСЯ: W, (2,3)} \\
            }
        }
    }
    & & \\ \\

    \theutterance \stepcounter{utterance}  
    & & & \multicolumn{2}{p{0.3\linewidth}}{
        \cellcolor[rgb]{0.9,0.9,0.9}{
            \makecell[{{p{\linewidth}}}]{
                \texttt{\tiny{[GM$|$GM]}}
                \texttt{W moved to (2, 3)} \\
            }
        }
    }
    & & \\ \\

    \theutterance \stepcounter{utterance}  
    & & \multicolumn{4}{p{0.6\linewidth}}{
        \cellcolor[rgb]{0.9,0.9,0.9}{
            \makecell[{{p{\linewidth}}}]{
                \texttt{\tiny{[GM$\rangle$P2]}}
                \texttt{Успешно перемещён объект "C" в (2,6). Ваша обновлённая сетка выглядит так:} \\
\texttt{```} \\
\texttt{ 1234567} \\
\texttt{╔════╤══╗} \\
\texttt{║◌◌◌T│◌◌║ 1} \\
\texttt{║◌◌◌◌│◌◌║ 2} \\
\texttt{║◌◌◌◌│◌◌║ 3} \\
\texttt{║W◌◌◌│◌◌║ 4} \\
\texttt{╟──┐◌│◌◌║ 5} \\
\texttt{║HC├─┼──╢ 6} \\
\texttt{║I◌│◌│◌◌║ 7} \\
\texttt{╚══╧═╧══╝} \\
\texttt{```} \\
\texttt{Сейчас вы играете раунд 5 из максимум 20.} \\
\texttt{Вы вместе накопили 4 из 12 штрафов.} \\
\texttt{Другой игрок переместил объект на своей сетке.} \\
\\ 
\\ 
\texttt{Какая твоя следующая команда?} \\
            }
        }
    }
    & \\ \\

    \theutterance \stepcounter{utterance}  
    & & & \multicolumn{4}{p{0.6\linewidth}}{
        \cellcolor[rgb]{0.9,0.9,0.9}{
            \makecell[{{p{\linewidth}}}]{
                \texttt{\tiny{[GM$\langle$P2]}}
                \texttt{ДВИГАТЬСЯ: T, (2, 2)} \\
            }
        }
    }
     \\ \\

    \theutterance \stepcounter{utterance}  
    & & & \multicolumn{2}{p{0.3\linewidth}}{
        \cellcolor[rgb]{0.9,0.9,0.9}{
            \makecell[{{p{\linewidth}}}]{
                \texttt{\tiny{[GM$|$GM]}}
                \texttt{ДВИГАТЬСЯ: T, (2, 2)} \\
            }
        }
    }
    & & \\ \\

    \theutterance \stepcounter{utterance}  
    & & & \multicolumn{2}{p{0.3\linewidth}}{
        \cellcolor[rgb]{0.9,0.9,0.9}{
            \makecell[{{p{\linewidth}}}]{
                \texttt{\tiny{[GM$|$GM]}}
                \texttt{T moved to (2, 2)} \\
            }
        }
    }
    & & \\ \\

    \theutterance \stepcounter{utterance}  
    & & \multicolumn{4}{p{0.6\linewidth}}{
        \cellcolor[rgb]{0.9,0.9,0.9}{
            \makecell[{{p{\linewidth}}}]{
                \texttt{\tiny{[P1$\langle$GM]}}
                \texttt{Успешно перемещён объект "W" в (2,3). Ваша обновлённая сетка выглядит так:} \\
\texttt{```} \\
\texttt{ 1234567} \\
\texttt{╔════╤══╗} \\
\texttt{║◌◌◌◌│◌◌║ 1} \\
\texttt{║◌◌◌◌│◌◌║ 2} \\
\texttt{║◌W◌◌│◌◌║ 3} \\
\texttt{║TH◌◌│I◌║ 4} \\
\texttt{╟──┐◌│◌◌║ 5} \\
\texttt{║◌C├─┼──╢ 6} \\
\texttt{║◌◌│◌│◌◌║ 7} \\
\texttt{╚══╧═╧══╝} \\
\texttt{```} \\
\texttt{Сейчас вы играете раунд 5 из максимум 20.} \\
\texttt{Вы вместе накопили 4 из 12 штрафов.} \\
\texttt{Другой игрок переместил объект на своей сетке.} \\
\\ 
\\ 
\texttt{Какая твоя следующая команда?} \\
            }
        }
    }
    & \\ \\

    \theutterance \stepcounter{utterance}  
    & \multicolumn{4}{p{0.6\linewidth}}{
        \cellcolor[rgb]{0.9,0.9,0.9}{
            \makecell[{{p{\linewidth}}}]{
                \texttt{\tiny{[P1$\rangle$GM]}}
                \texttt{ДВИГАТЬСЯ: T, (2,2)} \\
            }
        }
    }
    & & \\ \\

    \theutterance \stepcounter{utterance}  
    & & & \multicolumn{2}{p{0.3\linewidth}}{
        \cellcolor[rgb]{0.9,0.9,0.9}{
            \makecell[{{p{\linewidth}}}]{
                \texttt{\tiny{[GM$|$GM]}}
                \texttt{ДВИГАТЬСЯ: T, (2,2)} \\
            }
        }
    }
    & & \\ \\

    \theutterance \stepcounter{utterance}  
    & & & \multicolumn{2}{p{0.3\linewidth}}{
        \cellcolor[rgb]{0.9,0.9,0.9}{
            \makecell[{{p{\linewidth}}}]{
                \texttt{\tiny{[GM$|$GM]}}
                \texttt{T moved to (2, 2)} \\
            }
        }
    }
    & & \\ \\

    \theutterance \stepcounter{utterance}  
    & & \multicolumn{4}{p{0.6\linewidth}}{
        \cellcolor[rgb]{0.9,0.9,0.9}{
            \makecell[{{p{\linewidth}}}]{
                \texttt{\tiny{[GM$\rangle$P2]}}
                \texttt{Успешно перемещён объект "T" в (2,2). Ваша обновлённая сетка выглядит так:} \\
\texttt{```} \\
\texttt{ 1234567} \\
\texttt{╔════╤══╗} \\
\texttt{║◌◌◌◌│◌◌║ 1} \\
\texttt{║◌T◌◌│◌◌║ 2} \\
\texttt{║◌◌◌◌│◌◌║ 3} \\
\texttt{║W◌◌◌│◌◌║ 4} \\
\texttt{╟──┐◌│◌◌║ 5} \\
\texttt{║HC├─┼──╢ 6} \\
\texttt{║I◌│◌│◌◌║ 7} \\
\texttt{╚══╧═╧══╝} \\
\texttt{```} \\
\texttt{Сейчас вы играете раунд 6 из максимум 20.} \\
\texttt{Вы вместе накопили 4 из 12 штрафов.} \\
\texttt{Другой игрок переместил объект на своей сетке.} \\
\\ 
\\ 
\texttt{Какая твоя следующая команда?} \\
            }
        }
    }
    & \\ \\

    \theutterance \stepcounter{utterance}  
    & & & \multicolumn{4}{p{0.6\linewidth}}{
        \cellcolor[rgb]{0.9,0.9,0.9}{
            \makecell[{{p{\linewidth}}}]{
                \texttt{\tiny{[GM$\langle$P2]}}
                \texttt{ДВИГАТЬСЯ: W, (2, 3)} \\
            }
        }
    }
     \\ \\

    \theutterance \stepcounter{utterance}  
    & & & \multicolumn{2}{p{0.3\linewidth}}{
        \cellcolor[rgb]{0.9,0.9,0.9}{
            \makecell[{{p{\linewidth}}}]{
                \texttt{\tiny{[GM$|$GM]}}
                \texttt{ДВИГАТЬСЯ: W, (2, 3)} \\
            }
        }
    }
    & & \\ \\

    \theutterance \stepcounter{utterance}  
    & & & \multicolumn{2}{p{0.3\linewidth}}{
        \cellcolor[rgb]{0.9,0.9,0.9}{
            \makecell[{{p{\linewidth}}}]{
                \texttt{\tiny{[GM$|$GM]}}
                \texttt{W moved to (2, 3)} \\
            }
        }
    }
    & & \\ \\

    \theutterance \stepcounter{utterance}  
    & & \multicolumn{4}{p{0.6\linewidth}}{
        \cellcolor[rgb]{0.9,0.9,0.9}{
            \makecell[{{p{\linewidth}}}]{
                \texttt{\tiny{[P1$\langle$GM]}}
                \texttt{Успешно перемещён объект "T" в (2,2). Ваша обновлённая сетка выглядит так:} \\
\texttt{```} \\
\texttt{ 1234567} \\
\texttt{╔════╤══╗} \\
\texttt{║◌◌◌◌│◌◌║ 1} \\
\texttt{║◌T◌◌│◌◌║ 2} \\
\texttt{║◌W◌◌│◌◌║ 3} \\
\texttt{║◌H◌◌│I◌║ 4} \\
\texttt{╟──┐◌│◌◌║ 5} \\
\texttt{║◌C├─┼──╢ 6} \\
\texttt{║◌◌│◌│◌◌║ 7} \\
\texttt{╚══╧═╧══╝} \\
\texttt{```} \\
\texttt{Сейчас вы играете раунд 6 из максимум 20.} \\
\texttt{Вы вместе накопили 4 из 12 штрафов.} \\
\texttt{Другой игрок переместил объект на своей сетке.} \\
\\ 
\\ 
\texttt{Какая твоя следующая команда?} \\
            }
        }
    }
    & \\ \\

    \theutterance \stepcounter{utterance}  
    & \multicolumn{4}{p{0.6\linewidth}}{
        \cellcolor[rgb]{0.9,0.9,0.9}{
            \makecell[{{p{\linewidth}}}]{
                \texttt{\tiny{[P1$\rangle$GM]}}
                \texttt{ДВИГАТЬСЯ: I, (2,5)} \\
            }
        }
    }
    & & \\ \\

    \theutterance \stepcounter{utterance}  
    & & & \multicolumn{2}{p{0.3\linewidth}}{
        \cellcolor[rgb]{0.9,0.9,0.9}{
            \makecell[{{p{\linewidth}}}]{
                \texttt{\tiny{[GM$|$GM]}}
                \texttt{ДВИГАТЬСЯ: I, (2,5)} \\
            }
        }
    }
    & & \\ \\

    \theutterance \stepcounter{utterance}  
    & & & \multicolumn{2}{p{0.3\linewidth}}{
        \cellcolor[rgb]{0.9,0.9,0.9}{
            \makecell[{{p{\linewidth}}}]{
                \texttt{\tiny{[GM$|$GM]}}
                \texttt{Штраф: (2,5) не пусто, а содержит "─".} \\
\\ 
\texttt{Вы вместе накопили 5 из 12 штрафов.} \\
\texttt{Пожалуйста, попробуйте ещё раз!} \\
            }
        }
    }
    & & \\ \\

    \theutterance \stepcounter{utterance}  
    & & \multicolumn{4}{p{0.6\linewidth}}{
        \cellcolor[rgb]{0.9,0.9,0.9}{
            \makecell[{{p{\linewidth}}}]{
                \texttt{\tiny{[P1$\langle$GM]}}
                \texttt{Штраф: (2,5) не пусто, а содержит "─".} \\
\\ 
\texttt{Вы вместе накопили 5 из 12 штрафов.} \\
\texttt{Пожалуйста, попробуйте ещё раз!} \\
            }
        }
    }
    & \\ \\

    \theutterance \stepcounter{utterance}  
    & \multicolumn{4}{p{0.6\linewidth}}{
        \cellcolor[rgb]{0.9,0.9,0.9}{
            \makecell[{{p{\linewidth}}}]{
                \texttt{\tiny{[P1$\rangle$GM]}}
                \texttt{ДВИГАТЬСЯ: I, (5,5)} \\
            }
        }
    }
    & & \\ \\

    \theutterance \stepcounter{utterance}  
    & & & \multicolumn{2}{p{0.3\linewidth}}{
        \cellcolor[rgb]{0.9,0.9,0.9}{
            \makecell[{{p{\linewidth}}}]{
                \texttt{\tiny{[GM$|$GM]}}
                \texttt{ДВИГАТЬСЯ: I, (5,5)} \\
            }
        }
    }
    & & \\ \\

    \theutterance \stepcounter{utterance}  
    & & & \multicolumn{2}{p{0.3\linewidth}}{
        \cellcolor[rgb]{0.9,0.9,0.9}{
            \makecell[{{p{\linewidth}}}]{
                \texttt{\tiny{[GM$|$GM]}}
                \texttt{Штраф: (5,5) не пусто, а содержит "│".} \\
\\ 
\texttt{Вы вместе накопили 6 из 12 штрафов.} \\
\texttt{Пожалуйста, попробуйте ещё раз!} \\
            }
        }
    }
    & & \\ \\

    \theutterance \stepcounter{utterance}  
    & & \multicolumn{4}{p{0.6\linewidth}}{
        \cellcolor[rgb]{0.9,0.9,0.9}{
            \makecell[{{p{\linewidth}}}]{
                \texttt{\tiny{[P1$\langle$GM]}}
                \texttt{Штраф: (5,5) не пусто, а содержит "│".} \\
\\ 
\texttt{Вы вместе накопили 6 из 12 штрафов.} \\
\texttt{Пожалуйста, попробуйте ещё раз!} \\
            }
        }
    }
    & \\ \\

    \theutterance \stepcounter{utterance}  
    & \multicolumn{4}{p{0.6\linewidth}}{
        \cellcolor[rgb]{0.9,0.9,0.9}{
            \makecell[{{p{\linewidth}}}]{
                \texttt{\tiny{[P1$\rangle$GM]}}
                \texttt{ДВИГАТЬСЯ: I, (5,1)} \\
            }
        }
    }
    & & \\ \\

    \theutterance \stepcounter{utterance}  
    & & & \multicolumn{2}{p{0.3\linewidth}}{
        \cellcolor[rgb]{0.9,0.9,0.9}{
            \makecell[{{p{\linewidth}}}]{
                \texttt{\tiny{[GM$|$GM]}}
                \texttt{ДВИГАТЬСЯ: I, (5,1)} \\
            }
        }
    }
    & & \\ \\

    \theutterance \stepcounter{utterance}  
    & & & \multicolumn{2}{p{0.3\linewidth}}{
        \cellcolor[rgb]{0.9,0.9,0.9}{
            \makecell[{{p{\linewidth}}}]{
                \texttt{\tiny{[GM$|$GM]}}
                \texttt{Штраф: (5,1) не пусто, а содержит "│".} \\
\\ 
\texttt{Вы вместе накопили 7 из 12 штрафов.} \\
\texttt{Пожалуйста, попробуйте ещё раз!} \\
            }
        }
    }
    & & \\ \\

    \theutterance \stepcounter{utterance}  
    & & \multicolumn{4}{p{0.6\linewidth}}{
        \cellcolor[rgb]{0.9,0.9,0.9}{
            \makecell[{{p{\linewidth}}}]{
                \texttt{\tiny{[P1$\langle$GM]}}
                \texttt{Штраф: (5,1) не пусто, а содержит "│".} \\
\\ 
\texttt{Вы вместе накопили 7 из 12 штрафов.} \\
\texttt{Пожалуйста, попробуйте ещё раз!} \\
            }
        }
    }
    & \\ \\

    \theutterance \stepcounter{utterance}  
    & \multicolumn{4}{p{0.6\linewidth}}{
        \cellcolor[rgb]{0.9,0.9,0.9}{
            \makecell[{{p{\linewidth}}}]{
                \texttt{\tiny{[P1$\rangle$GM]}}
                \texttt{ДВИГАТЬСЯ: I, (6,1)} \\
            }
        }
    }
    & & \\ \\

    \theutterance \stepcounter{utterance}  
    & & & \multicolumn{2}{p{0.3\linewidth}}{
        \cellcolor[rgb]{0.9,0.9,0.9}{
            \makecell[{{p{\linewidth}}}]{
                \texttt{\tiny{[GM$|$GM]}}
                \texttt{ДВИГАТЬСЯ: I, (6,1)} \\
            }
        }
    }
    & & \\ \\

    \theutterance \stepcounter{utterance}  
    & & & \multicolumn{2}{p{0.3\linewidth}}{
        \cellcolor[rgb]{0.9,0.9,0.9}{
            \makecell[{{p{\linewidth}}}]{
                \texttt{\tiny{[GM$|$GM]}}
                \texttt{I moved to (6, 1)} \\
            }
        }
    }
    & & \\ \\

    \theutterance \stepcounter{utterance}  
    & & \multicolumn{4}{p{0.6\linewidth}}{
        \cellcolor[rgb]{0.9,0.9,0.9}{
            \makecell[{{p{\linewidth}}}]{
                \texttt{\tiny{[GM$\rangle$P2]}}
                \texttt{Успешно перемещён объект "W" в (2,3). Ваша обновлённая сетка выглядит так:} \\
\texttt{```} \\
\texttt{ 1234567} \\
\texttt{╔════╤══╗} \\
\texttt{║◌◌◌◌│◌◌║ 1} \\
\texttt{║◌T◌◌│◌◌║ 2} \\
\texttt{║◌W◌◌│◌◌║ 3} \\
\texttt{║◌◌◌◌│◌◌║ 4} \\
\texttt{╟──┐◌│◌◌║ 5} \\
\texttt{║HC├─┼──╢ 6} \\
\texttt{║I◌│◌│◌◌║ 7} \\
\texttt{╚══╧═╧══╝} \\
\texttt{```} \\
\texttt{Сейчас вы играете раунд 7 из максимум 20.} \\
\texttt{Вы вместе накопили 7 из 12 штрафов.} \\
\texttt{Другой игрок переместил объект на своей сетке.} \\
\\ 
\\ 
\texttt{Какая твоя следующая команда?} \\
            }
        }
    }
    & \\ \\

    \theutterance \stepcounter{utterance}  
    & & & \multicolumn{4}{p{0.6\linewidth}}{
        \cellcolor[rgb]{0.9,0.9,0.9}{
            \makecell[{{p{\linewidth}}}]{
                \texttt{\tiny{[GM$\langle$P2]}}
                \texttt{ДВИГАТЬСЯ: H, (2, 4)} \\
            }
        }
    }
     \\ \\

    \theutterance \stepcounter{utterance}  
    & & & \multicolumn{2}{p{0.3\linewidth}}{
        \cellcolor[rgb]{0.9,0.9,0.9}{
            \makecell[{{p{\linewidth}}}]{
                \texttt{\tiny{[GM$|$GM]}}
                \texttt{ДВИГАТЬСЯ: H, (2, 4)} \\
            }
        }
    }
    & & \\ \\

    \theutterance \stepcounter{utterance}  
    & & & \multicolumn{2}{p{0.3\linewidth}}{
        \cellcolor[rgb]{0.9,0.9,0.9}{
            \makecell[{{p{\linewidth}}}]{
                \texttt{\tiny{[GM$|$GM]}}
                \texttt{H moved to (2, 4)} \\
            }
        }
    }
    & & \\ \\

    \theutterance \stepcounter{utterance}  
    & & \multicolumn{4}{p{0.6\linewidth}}{
        \cellcolor[rgb]{0.9,0.9,0.9}{
            \makecell[{{p{\linewidth}}}]{
                \texttt{\tiny{[P1$\langle$GM]}}
                \texttt{Успешно перемещён объект "I" в (6,1). Ваша обновлённая сетка выглядит так:} \\
\texttt{```} \\
\texttt{ 1234567} \\
\texttt{╔════╤══╗} \\
\texttt{║◌◌◌◌│I◌║ 1} \\
\texttt{║◌T◌◌│◌◌║ 2} \\
\texttt{║◌W◌◌│◌◌║ 3} \\
\texttt{║◌H◌◌│◌◌║ 4} \\
\texttt{╟──┐◌│◌◌║ 5} \\
\texttt{║◌C├─┼──╢ 6} \\
\texttt{║◌◌│◌│◌◌║ 7} \\
\texttt{╚══╧═╧══╝} \\
\texttt{```} \\
\texttt{Сейчас вы играете раунд 7 из максимум 20.} \\
\texttt{Вы вместе накопили 7 из 12 штрафов.} \\
\texttt{Другой игрок переместил объект на своей сетке.} \\
\\ 
\\ 
\texttt{Какая твоя следующая команда?} \\
            }
        }
    }
    & \\ \\

    \theutterance \stepcounter{utterance}  
    & \multicolumn{4}{p{0.6\linewidth}}{
        \cellcolor[rgb]{0.9,0.9,0.9}{
            \makecell[{{p{\linewidth}}}]{
                \texttt{\tiny{[P1$\rangle$GM]}}
                \texttt{ДВИГАТЬСЯ: I, (2,5)} \\
            }
        }
    }
    & & \\ \\

    \theutterance \stepcounter{utterance}  
    & & & \multicolumn{2}{p{0.3\linewidth}}{
        \cellcolor[rgb]{0.9,0.9,0.9}{
            \makecell[{{p{\linewidth}}}]{
                \texttt{\tiny{[GM$|$GM]}}
                \texttt{ДВИГАТЬСЯ: I, (2,5)} \\
            }
        }
    }
    & & \\ \\

    \theutterance \stepcounter{utterance}  
    & & & \multicolumn{2}{p{0.3\linewidth}}{
        \cellcolor[rgb]{0.9,0.9,0.9}{
            \makecell[{{p{\linewidth}}}]{
                \texttt{\tiny{[GM$|$GM]}}
                \texttt{Штраф: (2,5) не пусто, а содержит "─".} \\
\\ 
\texttt{Вы вместе накопили 8 из 12 штрафов.} \\
\texttt{Пожалуйста, попробуйте ещё раз!} \\
            }
        }
    }
    & & \\ \\

    \theutterance \stepcounter{utterance}  
    & & \multicolumn{4}{p{0.6\linewidth}}{
        \cellcolor[rgb]{0.9,0.9,0.9}{
            \makecell[{{p{\linewidth}}}]{
                \texttt{\tiny{[P1$\langle$GM]}}
                \texttt{Штраф: (2,5) не пусто, а содержит "─".} \\
\\ 
\texttt{Вы вместе накопили 8 из 12 штрафов.} \\
\texttt{Пожалуйста, попробуйте ещё раз!} \\
            }
        }
    }
    & \\ \\

    \theutterance \stepcounter{utterance}  
    & \multicolumn{4}{p{0.6\linewidth}}{
        \cellcolor[rgb]{0.9,0.9,0.9}{
            \makecell[{{p{\linewidth}}}]{
                \texttt{\tiny{[P1$\rangle$GM]}}
                \texttt{ДВИГАТЬСЯ: I, (5,4)} \\
            }
        }
    }
    & & \\ \\

    \theutterance \stepcounter{utterance}  
    & & & \multicolumn{2}{p{0.3\linewidth}}{
        \cellcolor[rgb]{0.9,0.9,0.9}{
            \makecell[{{p{\linewidth}}}]{
                \texttt{\tiny{[GM$|$GM]}}
                \texttt{ДВИГАТЬСЯ: I, (5,4)} \\
            }
        }
    }
    & & \\ \\

    \theutterance \stepcounter{utterance}  
    & & & \multicolumn{2}{p{0.3\linewidth}}{
        \cellcolor[rgb]{0.9,0.9,0.9}{
            \makecell[{{p{\linewidth}}}]{
                \texttt{\tiny{[GM$|$GM]}}
                \texttt{Штраф: (5,4) не пусто, а содержит "│".} \\
\\ 
\texttt{Вы вместе накопили 9 из 12 штрафов.} \\
\texttt{Пожалуйста, попробуйте ещё раз!} \\
            }
        }
    }
    & & \\ \\

    \theutterance \stepcounter{utterance}  
    & & \multicolumn{4}{p{0.6\linewidth}}{
        \cellcolor[rgb]{0.9,0.9,0.9}{
            \makecell[{{p{\linewidth}}}]{
                \texttt{\tiny{[P1$\langle$GM]}}
                \texttt{Штраф: (5,4) не пусто, а содержит "│".} \\
\\ 
\texttt{Вы вместе накопили 9 из 12 штрафов.} \\
\texttt{Пожалуйста, попробуйте ещё раз!} \\
            }
        }
    }
    & \\ \\

    \theutterance \stepcounter{utterance}  
    & \multicolumn{4}{p{0.6\linewidth}}{
        \cellcolor[rgb]{0.9,0.9,0.9}{
            \makecell[{{p{\linewidth}}}]{
                \texttt{\tiny{[P1$\rangle$GM]}}
                \texttt{ДВИГАТЬСЯ: I, (6,4)} \\
            }
        }
    }
    & & \\ \\

    \theutterance \stepcounter{utterance}  
    & & & \multicolumn{2}{p{0.3\linewidth}}{
        \cellcolor[rgb]{0.9,0.9,0.9}{
            \makecell[{{p{\linewidth}}}]{
                \texttt{\tiny{[GM$|$GM]}}
                \texttt{ДВИГАТЬСЯ: I, (6,4)} \\
            }
        }
    }
    & & \\ \\

    \theutterance \stepcounter{utterance}  
    & & & \multicolumn{2}{p{0.3\linewidth}}{
        \cellcolor[rgb]{0.9,0.9,0.9}{
            \makecell[{{p{\linewidth}}}]{
                \texttt{\tiny{[GM$|$GM]}}
                \texttt{I moved to (6, 4)} \\
            }
        }
    }
    & & \\ \\

    \theutterance \stepcounter{utterance}  
    & & \multicolumn{4}{p{0.6\linewidth}}{
        \cellcolor[rgb]{0.9,0.9,0.9}{
            \makecell[{{p{\linewidth}}}]{
                \texttt{\tiny{[GM$\rangle$P2]}}
                \texttt{Успешно перемещён объект "H" в (2,4). Ваша обновлённая сетка выглядит так:} \\
\texttt{```} \\
\texttt{ 1234567} \\
\texttt{╔════╤══╗} \\
\texttt{║◌◌◌◌│◌◌║ 1} \\
\texttt{║◌T◌◌│◌◌║ 2} \\
\texttt{║◌W◌◌│◌◌║ 3} \\
\texttt{║◌H◌◌│◌◌║ 4} \\
\texttt{╟──┐◌│◌◌║ 5} \\
\texttt{║◌C├─┼──╢ 6} \\
\texttt{║I◌│◌│◌◌║ 7} \\
\texttt{╚══╧═╧══╝} \\
\texttt{```} \\
\texttt{Сейчас вы играете раунд 8 из максимум 20.} \\
\texttt{Вы вместе накопили 9 из 12 штрафов.} \\
\texttt{Другой игрок переместил объект на своей сетке.} \\
\\ 
\\ 
\texttt{Какая твоя следующая команда?} \\
            }
        }
    }
    & \\ \\

    \theutterance \stepcounter{utterance}  
    & & & \multicolumn{4}{p{0.6\linewidth}}{
        \cellcolor[rgb]{0.9,0.9,0.9}{
            \makecell[{{p{\linewidth}}}]{
                \texttt{\tiny{[GM$\langle$P2]}}
                \texttt{ДВИГАТЬСЯ: I, (2, 5)} \\
            }
        }
    }
     \\ \\

    \theutterance \stepcounter{utterance}  
    & & & \multicolumn{2}{p{0.3\linewidth}}{
        \cellcolor[rgb]{0.9,0.9,0.9}{
            \makecell[{{p{\linewidth}}}]{
                \texttt{\tiny{[GM$|$GM]}}
                \texttt{ДВИГАТЬСЯ: I, (2, 5)} \\
            }
        }
    }
    & & \\ \\

    \theutterance \stepcounter{utterance}  
    & & & \multicolumn{2}{p{0.3\linewidth}}{
        \cellcolor[rgb]{0.9,0.9,0.9}{
            \makecell[{{p{\linewidth}}}]{
                \texttt{\tiny{[GM$|$GM]}}
                \texttt{Штраф: (2,5) не пусто, а содержит "─".} \\
\\ 
\texttt{Вы вместе накопили 10 из 12 штрафов.} \\
\texttt{Пожалуйста, попробуйте ещё раз!} \\
            }
        }
    }
    & & \\ \\

    \theutterance \stepcounter{utterance}  
    & & \multicolumn{4}{p{0.6\linewidth}}{
        \cellcolor[rgb]{0.9,0.9,0.9}{
            \makecell[{{p{\linewidth}}}]{
                \texttt{\tiny{[GM$\rangle$P2]}}
                \texttt{Штраф: (2,5) не пусто, а содержит "─".} \\
\\ 
\texttt{Вы вместе накопили 10 из 12 штрафов.} \\
\texttt{Пожалуйста, попробуйте ещё раз!} \\
            }
        }
    }
    & \\ \\

    \theutterance \stepcounter{utterance}  
    & & & \multicolumn{4}{p{0.6\linewidth}}{
        \cellcolor[rgb]{0.9,0.9,0.9}{
            \makecell[{{p{\linewidth}}}]{
                \texttt{\tiny{[GM$\langle$P2]}}
                \texttt{ДВИГАТЬСЯ: I, (1, 5)} \\
            }
        }
    }
     \\ \\

    \theutterance \stepcounter{utterance}  
    & & & \multicolumn{2}{p{0.3\linewidth}}{
        \cellcolor[rgb]{0.9,0.9,0.9}{
            \makecell[{{p{\linewidth}}}]{
                \texttt{\tiny{[GM$|$GM]}}
                \texttt{ДВИГАТЬСЯ: I, (1, 5)} \\
            }
        }
    }
    & & \\ \\

    \theutterance \stepcounter{utterance}  
    & & & \multicolumn{2}{p{0.3\linewidth}}{
        \cellcolor[rgb]{0.9,0.9,0.9}{
            \makecell[{{p{\linewidth}}}]{
                \texttt{\tiny{[GM$|$GM]}}
                \texttt{Штраф: (1,5) не пусто, а содержит "─".} \\
\\ 
\texttt{Вы вместе накопили 11 из 12 штрафов.} \\
\texttt{Пожалуйста, попробуйте ещё раз!} \\
            }
        }
    }
    & & \\ \\

    \theutterance \stepcounter{utterance}  
    & & \multicolumn{4}{p{0.6\linewidth}}{
        \cellcolor[rgb]{0.9,0.9,0.9}{
            \makecell[{{p{\linewidth}}}]{
                \texttt{\tiny{[GM$\rangle$P2]}}
                \texttt{Штраф: (1,5) не пусто, а содержит "─".} \\
\\ 
\texttt{Вы вместе накопили 11 из 12 штрафов.} \\
\texttt{Пожалуйста, попробуйте ещё раз!} \\
            }
        }
    }
    & \\ \\

    \theutterance \stepcounter{utterance}  
    & & & \multicolumn{4}{p{0.6\linewidth}}{
        \cellcolor[rgb]{0.9,0.9,0.9}{
            \makecell[{{p{\linewidth}}}]{
                \texttt{\tiny{[GM$\langle$P2]}}
                \texttt{ДВИГАТЬСЯ: I, (2, 7)} \\
            }
        }
    }
     \\ \\

    \theutterance \stepcounter{utterance}  
    & & & \multicolumn{2}{p{0.3\linewidth}}{
        \cellcolor[rgb]{0.9,0.9,0.9}{
            \makecell[{{p{\linewidth}}}]{
                \texttt{\tiny{[GM$|$GM]}}
                \texttt{ДВИГАТЬСЯ: I, (2, 7)} \\
            }
        }
    }
    & & \\ \\

    \theutterance \stepcounter{utterance}  
    & & & \multicolumn{2}{p{0.3\linewidth}}{
        \cellcolor[rgb]{0.9,0.9,0.9}{
            \makecell[{{p{\linewidth}}}]{
                \texttt{\tiny{[GM$|$GM]}}
                \texttt{I moved to (2, 7)} \\
            }
        }
    }
    & & \\ \\

    \theutterance \stepcounter{utterance}  
    & & \multicolumn{4}{p{0.6\linewidth}}{
        \cellcolor[rgb]{0.9,0.9,0.9}{
            \makecell[{{p{\linewidth}}}]{
                \texttt{\tiny{[P1$\langle$GM]}}
                \texttt{Успешно перемещён объект "I" в (6,4). Ваша обновлённая сетка выглядит так:} \\
\texttt{```} \\
\texttt{ 1234567} \\
\texttt{╔════╤══╗} \\
\texttt{║◌◌◌◌│◌◌║ 1} \\
\texttt{║◌T◌◌│◌◌║ 2} \\
\texttt{║◌W◌◌│◌◌║ 3} \\
\texttt{║◌H◌◌│I◌║ 4} \\
\texttt{╟──┐◌│◌◌║ 5} \\
\texttt{║◌C├─┼──╢ 6} \\
\texttt{║◌◌│◌│◌◌║ 7} \\
\texttt{╚══╧═╧══╝} \\
\texttt{```} \\
\texttt{Сейчас вы играете раунд 8 из максимум 20.} \\
\texttt{Вы вместе накопили 11 из 12 штрафов.} \\
\texttt{Другой игрок переместил объект на своей сетке.} \\
\\ 
\\ 
\texttt{Какая твоя следующая команда?} \\
            }
        }
    }
    & \\ \\

    \theutterance \stepcounter{utterance}  
    & \multicolumn{4}{p{0.6\linewidth}}{
        \cellcolor[rgb]{0.9,0.9,0.9}{
            \makecell[{{p{\linewidth}}}]{
                \texttt{\tiny{[P1$\rangle$GM]}}
                \texttt{ДВИГАТЬСЯ: I, (2,5)} \\
\\ 
\texttt{(Примечание: Я вижу, что (2,5) содержит символ '─' в сетке, но в правилах сказано, что я должен перемещать в пустую клетку '◌'. Посмотрев на сетку внимательнее, я вижу, что (2,5) — это часть разделительной линии. Попробую найти свободную клетку рядом с целью. Цель для I была (2,5). Попробую переместить I в (3,5), если это возможно, но (3,5) тоже часть линии. Попробую (2,7) или (1,5)? Нет, (2,5) — это клетка под (2,4). В сетке (2,4) — это 'H', (2,5) — это '◌' под 'H'?)} \\
\\ 
\texttt{*Перепроверяю сетку:*} \\
\texttt{(2,1) ◌} \\
\texttt{(2,2) T} \\
\texttt{(2,3) W} \\
\texttt{(2,4) H} \\
\texttt{(2,5) ◌ (в строке 5, под H)} \\
\texttt{(2,6) C} \\
\\ 
\texttt{Похоже, в предыдущем ходу я ошибся, считая (2,5) занятым линией. В строке 5: `╟──┐◌│◌◌║ 5`. Координата (2,5) — это символ '─'. Координата (4,5) — это '◌'.} \\
\texttt{Попробую переместить I в (2,5) еще раз, если это была ошибка интерпретации, или выберу безопасную клетку.} \\
\\ 
\texttt{Стоп, (2,5) — это точно '─'.} \\
\texttt{Попробую:} \\
\texttt{ДВИГАТЬСЯ: I, (2,7)} \\
            }
        }
    }
    & & \\ \\

    \theutterance \stepcounter{utterance}  
    & & & \multicolumn{2}{p{0.3\linewidth}}{
        \cellcolor[rgb]{0.9,0.9,0.9}{
            \makecell[{{p{\linewidth}}}]{
                \texttt{\tiny{[GM$|$GM]}}
                \texttt{ДВИГАТЬСЯ: I, (2,5)} \\
\\ 
\texttt{(Примечание: Я вижу, что (2,5) содержит символ '─' в сетке, но в правилах сказано, что я должен перемещать в пустую клетку '◌'. Посмотрев на сетку внимательнее, я вижу, что (2,5) — это часть разделительной линии. Попробую найти свободную клетку рядом с целью. Цель для I была (2,5). Попробую переместить I в (3,5), если это возможно, но (3,5) тоже часть линии. Попробую (2,7) или (1,5)? Нет, (2,5) — это клетка под (2,4). В сетке (2,4) — это 'H', (2,5) — это '◌' под 'H'?)} \\
\\ 
\texttt{*Перепроверяю сетку:*} \\
\texttt{(2,1) ◌} \\
\texttt{(2,2) T} \\
\texttt{(2,3) W} \\
\texttt{(2,4) H} \\
\texttt{(2,5) ◌ (в строке 5, под H)} \\
\texttt{(2,6) C} \\
\\ 
\texttt{Похоже, в предыдущем ходу я ошибся, считая (2,5) занятым линией. В строке 5: `╟──┐◌│◌◌║ 5`. Координата (2,5) — это символ '─'. Координата (4,5) — это '◌'.} \\
\texttt{Попробую переместить I в (2,5) еще раз, если это была ошибка интерпретации, или выберу безопасную клетку.} \\
\\ 
\texttt{Стоп, (2,5) — это точно '─'.} \\
\texttt{Попробую:} \\
\texttt{ДВИГАТЬСЯ: I, (2,7)} \\
            }
        }
    }
    & & \\ \\

    \theutterance \stepcounter{utterance}  
    & & & \multicolumn{2}{p{0.3\linewidth}}{
        \cellcolor[rgb]{0.9,0.9,0.9}{
            \makecell[{{p{\linewidth}}}]{
                \texttt{\tiny{[GM$|$GM]}}
                \texttt{Invalid format: head is not empty: 'ДВИГАТЬСЯ: I, (2,5)} \\
\\ 
\texttt{(Примечание: Я вижу, что (2,5) содержит символ '─' в сетке, но в правилах сказано, что я должен перемещать в пустую клетку '◌'. Посмотрев на сетку внимательнее, я вижу, что (2,5) — это часть разделительной линии. Попробую найти свободную клетку рядом с целью. Цель для I была (2,5). Попробую переместить I в (3,5), если это возможно, но (3,5) тоже часть линии. Попробую (2,7) или (1,5)? Нет, (2,5) — это клетка под (2,4). В сетке (2,4) — это 'H', (2,5) — это '◌' под 'H'?)} \\
\\ 
\texttt{*Перепроверяю сетку:*} \\
\texttt{(2,1) ◌} \\
\texttt{(2,2) T} \\
\texttt{(2,3) W} \\
\texttt{(2,4) H} \\
\texttt{(2,5) ◌ (в строке 5, под H)} \\
\texttt{(2,6) C} \\
\\ 
\texttt{Похоже, в предыдущем ходу я ошибся, считая (2,5) занятым линией. В строке 5: ╟──┐◌│◌◌║ 5. Координата (2,5) — это символ '─'. Координата (4,5) — это '◌'.} \\
\texttt{Попробую переместить I в (2,5) еще раз, если это была ошибка интерпретации, или выберу безопасную клетку.} \\
\\ 
\texttt{Стоп, (2,5) — это точно '─'.} \\
\texttt{Попробую:} \\
\texttt{'} \\
            }
        }
    }
    & & \\ \\

    \theutterance \stepcounter{utterance}  
    & & \multicolumn{4}{p{0.6\linewidth}}{
        \cellcolor[rgb]{0.9,0.9,0.9}{
            \makecell[{{p{\linewidth}}}]{
                \texttt{\tiny{[P1$\langle$GM]}}
                \texttt{Штраф: Ваше сообщение не должно содержать ничего перед командой!} \\
\\ 
\texttt{Убедитесь, что ваш ответ содержит только `СКАЗАТЬ: <СООБЩЕНИЕ>` или `ДВИГАТЬСЯ: <ОБЪЕКТ>, (<X>, <Y>)` и ничего больше!} \\
\\ 
\texttt{Вы вместе накопили 12 из 12 штрафов.} \\
\texttt{Пожалуйста, попробуйте ещё раз!} \\
            }
        }
    }
    & \\ \\

    \theutterance \stepcounter{utterance}  
    & \multicolumn{4}{p{0.6\linewidth}}{
        \cellcolor[rgb]{0.9,0.9,0.9}{
            \makecell[{{p{\linewidth}}}]{
                \texttt{\tiny{[P1$\rangle$GM]}}
                \texttt{ДВИГАТЬСЯ: I, (2,7)} \\
            }
        }
    }
    & & \\ \\

    \theutterance \stepcounter{utterance}  
    & & & \multicolumn{2}{p{0.3\linewidth}}{
        \cellcolor[rgb]{0.9,0.9,0.9}{
            \makecell[{{p{\linewidth}}}]{
                \texttt{\tiny{[GM$|$GM]}}
                \texttt{ДВИГАТЬСЯ: I, (2,7)} \\
            }
        }
    }
    & & \\ \\

    \theutterance \stepcounter{utterance}  
    & & & \multicolumn{2}{p{0.3\linewidth}}{
        \cellcolor[rgb]{0.9,0.9,0.9}{
            \makecell[{{p{\linewidth}}}]{
                \texttt{\tiny{[GM$|$GM]}}
                \texttt{I moved to (2, 7)} \\
            }
        }
    }
    & & \\ \\

    \theutterance \stepcounter{utterance}  
    & & \multicolumn{4}{p{0.6\linewidth}}{
        \cellcolor[rgb]{0.9,0.9,0.9}{
            \makecell[{{p{\linewidth}}}]{
                \texttt{\tiny{[GM$\rangle$P2]}}
                \texttt{Успешно перемещён объект "I" в (2,7). Ваша обновлённая сетка выглядит так:} \\
\texttt{```} \\
\texttt{ 1234567} \\
\texttt{╔════╤══╗} \\
\texttt{║◌◌◌◌│◌◌║ 1} \\
\texttt{║◌T◌◌│◌◌║ 2} \\
\texttt{║◌W◌◌│◌◌║ 3} \\
\texttt{║◌H◌◌│◌◌║ 4} \\
\texttt{╟──┐◌│◌◌║ 5} \\
\texttt{║◌C├─┼──╢ 6} \\
\texttt{║◌I│◌│◌◌║ 7} \\
\texttt{╚══╧═╧══╝} \\
\texttt{```} \\
\texttt{Сейчас вы играете раунд 9 из максимум 20.} \\
\texttt{Вы вместе накопили 12 из 12 штрафов.} \\
\texttt{Другой игрок переместил объект на своей сетке.} \\
\\ 
\\ 
\texttt{Какая твоя следующая команда?} \\
            }
        }
    }
    & \\ \\

    \theutterance \stepcounter{utterance}  
    & & & \multicolumn{4}{p{0.6\linewidth}}{
        \cellcolor[rgb]{0.9,0.9,0.9}{
            \makecell[{{p{\linewidth}}}]{
                \texttt{\tiny{[GM$\langle$P2]}}
                \texttt{СКАЗАТЬ: Мы на грани проигрыша из{-}за штрафов. Пожалуйста, не делай больше ошибок. Я переместил I в (2,7), так как (2,5) и (1,5) были заняты линиями. Давай закончим, если твои объекты на месте.} \\
            }
        }
    }
     \\ \\

    \theutterance \stepcounter{utterance}  
    & & & \multicolumn{2}{p{0.3\linewidth}}{
        \cellcolor[rgb]{0.9,0.9,0.9}{
            \makecell[{{p{\linewidth}}}]{
                \texttt{\tiny{[GM$|$GM]}}
                \texttt{СКАЗАТЬ: Мы на грани проигрыша из{-}за штрафов. Пожалуйста, не делай больше ошибок. Я переместил I в (2,7), так как (2,5) и (1,5) были заняты линиями. Давай закончим, если твои объекты на месте.} \\
            }
        }
    }
    & & \\ \\

    \theutterance \stepcounter{utterance}  
    & & \multicolumn{4}{p{0.6\linewidth}}{
        \cellcolor[rgb]{0.9,0.9,0.9}{
            \makecell[{{p{\linewidth}}}]{
                \texttt{\tiny{[P1$\langle$GM]}}
                \texttt{Успешно перемещён объект "I" в (2,7). Ваша обновлённая сетка выглядит так:} \\
\texttt{```} \\
\texttt{ 1234567} \\
\texttt{╔════╤══╗} \\
\texttt{║◌◌◌◌│◌◌║ 1} \\
\texttt{║◌T◌◌│◌◌║ 2} \\
\texttt{║◌W◌◌│◌◌║ 3} \\
\texttt{║◌H◌◌│◌◌║ 4} \\
\texttt{╟──┐◌│◌◌║ 5} \\
\texttt{║◌C├─┼──╢ 6} \\
\texttt{║◌I│◌│◌◌║ 7} \\
\texttt{╚══╧═╧══╝} \\
\texttt{```} \\
\texttt{Сейчас вы играете раунд 9 из максимум 20.} \\
\texttt{Вы вместе накопили 12 из 12 штрафов.} \\
\texttt{Другой игрок отправил это сообщение:} \\
\\ 
\\ 
\texttt{"Мы на грани проигрыша из{-}за штрафов. Пожалуйста, не делай больше ошибок. Я переместил I в (2,7), так как (2,5) и (1,5) были заняты линиями. Давай закончим, если твои объекты на месте."} \\
\\ 
\\ 
\texttt{Какая твоя следующая команда?} \\
            }
        }
    }
    & \\ \\

    \theutterance \stepcounter{utterance}  
    & \multicolumn{4}{p{0.6\linewidth}}{
        \cellcolor[rgb]{0.9,0.9,0.9}{
            \makecell[{{p{\linewidth}}}]{
                \texttt{\tiny{[P1$\rangle$GM]}}
                \texttt{СКАЗАТЬ: Мои объекты сейчас на: T(2,2), W(2,3), H(2,4), I(2,7), C(2,6). Твоя цель для I была (2,5), но там линия. Если ты тоже на (2,7) или готов закончить, то готово?} \\
            }
        }
    }
    & & \\ \\

    \theutterance \stepcounter{utterance}  
    & & & \multicolumn{2}{p{0.3\linewidth}}{
        \cellcolor[rgb]{0.9,0.9,0.9}{
            \makecell[{{p{\linewidth}}}]{
                \texttt{\tiny{[GM$|$GM]}}
                \texttt{СКАЗАТЬ: Мои объекты сейчас на: T(2,2), W(2,3), H(2,4), I(2,7), C(2,6). Твоя цель для I была (2,5), но там линия. Если ты тоже на (2,7) или готов закончить, то готово?} \\
            }
        }
    }
    & & \\ \\

    \theutterance \stepcounter{utterance}  
    & & \multicolumn{4}{p{0.6\linewidth}}{
        \cellcolor[rgb]{0.9,0.9,0.9}{
            \makecell[{{p{\linewidth}}}]{
                \texttt{\tiny{[GM$\rangle$P2]}}
                \texttt{Ваше сообщение было передано другому игроку.} \\
\\ 
\\ 
\texttt{Сейчас вы играете раунд 10 из максимум 20.} \\
\texttt{Вы вместе накопили 12 из 12 штрафов.} \\
\texttt{Другой игрок отправил это сообщение:} \\
\\ 
\\ 
\texttt{"Мои объекты сейчас на: T(2,2), W(2,3), H(2,4), I(2,7), C(2,6). Твоя цель для I была (2,5), но там линия. Если ты тоже на (2,7) или готов закончить, то готово?"} \\
\\ 
\\ 
\texttt{Какая твоя следующая команда?} \\
            }
        }
    }
    & \\ \\

    \theutterance \stepcounter{utterance}  
    & & & \multicolumn{4}{p{0.6\linewidth}}{
        \cellcolor[rgb]{0.9,0.9,0.9}{
            \makecell[{{p{\linewidth}}}]{
                \texttt{\tiny{[GM$\langle$P2]}}
                \texttt{СКАЗАТЬ: готово!} \\
            }
        }
    }
     \\ \\

    \theutterance \stepcounter{utterance}  
    & & & \multicolumn{2}{p{0.3\linewidth}}{
        \cellcolor[rgb]{0.9,0.9,0.9}{
            \makecell[{{p{\linewidth}}}]{
                \texttt{\tiny{[GM$|$GM]}}
                \texttt{СКАЗАТЬ: готово!} \\
            }
        }
    }
    & & \\ \\

    \theutterance \stepcounter{utterance}  
    & & & \multicolumn{2}{p{0.3\linewidth}}{
        \cellcolor[rgb]{0.9,0.9,0.9}{
            \makecell[{{p{\linewidth}}}]{
                \texttt{\tiny{[GM$|$GM]}}
                \texttt{true} \\
            }
        }
    }
    & & \\ \\

    \theutterance \stepcounter{utterance}  
    & & & \multicolumn{2}{p{0.3\linewidth}}{
        \cellcolor[rgb]{0.9,0.9,0.9}{
            \makecell[{{p{\linewidth}}}]{
                \texttt{\tiny{[GM$|$GM]}}
                \texttt{Initial states:} \\
\texttt{```} \\
\texttt{Player 1:} \\
\texttt{ 1234567} \\
\texttt{╔════╤══╗} \\
\texttt{║H◌◌◌│◌◌║ 1} \\
\texttt{║◌W◌◌│◌◌║ 2} \\
\texttt{║◌C◌◌│◌◌║ 3} \\
\texttt{║T◌◌◌│I◌║ 4} \\
\texttt{╟──┐◌│◌◌║ 5} \\
\texttt{║◌◌├─┼──╢ 6} \\
\texttt{║◌◌│◌│◌◌║ 7} \\
\texttt{╚══╧═╧══╝} \\
\\ 
\\ 
\texttt{Player 2:} \\
\texttt{ 1234567} \\
\texttt{╔════╤══╗} \\
\texttt{║◌◌◌T│◌◌║ 1} \\
\texttt{║◌◌◌◌│◌◌║ 2} \\
\texttt{║◌◌◌◌│◌◌║ 3} \\
\texttt{║W◌◌◌│◌◌║ 4} \\
\texttt{╟──┐◌│◌◌║ 5} \\
\texttt{║H◌├─┼──╢ 6} \\
\texttt{║I◌│◌│◌C║ 7} \\
\texttt{╚══╧═╧══╝} \\
\\ 
\texttt{```} \\
            }
        }
    }
    & & \\ \\

    \theutterance \stepcounter{utterance}  
    & & & \multicolumn{2}{p{0.3\linewidth}}{
        \cellcolor[rgb]{0.9,0.9,0.9}{
            \makecell[{{p{\linewidth}}}]{
                \texttt{\tiny{[GM$|$GM]}}
                \texttt{End states:} \\
\texttt{```} \\
\texttt{Player 1:} \\
\texttt{ 1234567} \\
\texttt{╔════╤══╗} \\
\texttt{║◌◌◌◌│◌◌║ 1} \\
\texttt{║◌T◌◌│◌◌║ 2} \\
\texttt{║◌W◌◌│◌◌║ 3} \\
\texttt{║◌H◌◌│◌◌║ 4} \\
\texttt{╟──┐◌│◌◌║ 5} \\
\texttt{║◌C├─┼──╢ 6} \\
\texttt{║◌I│◌│◌◌║ 7} \\
\texttt{╚══╧═╧══╝} \\
\\ 
\\ 
\texttt{Player 2:} \\
\texttt{ 1234567} \\
\texttt{╔════╤══╗} \\
\texttt{║◌◌◌◌│◌◌║ 1} \\
\texttt{║◌T◌◌│◌◌║ 2} \\
\texttt{║◌W◌◌│◌◌║ 3} \\
\texttt{║◌H◌◌│◌◌║ 4} \\
\texttt{╟──┐◌│◌◌║ 5} \\
\texttt{║◌C├─┼──╢ 6} \\
\texttt{║◌I│◌│◌◌║ 7} \\
\texttt{╚══╧═╧══╝} \\
\\ 
\texttt{```} \\
            }
        }
    }
    & & \\ \\

    \theutterance \stepcounter{utterance}  
    & & & \multicolumn{2}{p{0.3\linewidth}}{
        \cellcolor[rgb]{0.9,0.9,0.9}{
            \makecell[{{p{\linewidth}}}]{
                \texttt{\tiny{[GM$|$GM]}}
                \texttt{* success: True} \\
\texttt{* lose: False} \\
\texttt{* aborted: False} \\
\texttt{{-}{-}{-}{-}{-}{-}{-}} \\
\texttt{* Shifts: 7.00} \\
\texttt{* Max Shifts: 8.00} \\
\texttt{* Min Shifts: 4.00} \\
\texttt{* End Distance Sum: 0.00} \\
\texttt{* Init Distance Sum: 23.71} \\
\texttt{* Expected Distance Sum: 20.95} \\
\texttt{* Penalties: 12.00} \\
\texttt{* Max Penalties: 12.00} \\
\texttt{* Rounds: 10.00} \\
\texttt{* Max Rounds: 20.00} \\
\texttt{* Object Count: 5.00} \\
            }
        }
    }
    & & \\ \\

\end{supertabular}
}

\end{document}
